\thispagestyle{plain}
\phantomsection
\addcontentsline{toc}{chapter}{Abstract}
\centerline{\zihao{-3}\heiti Abstract}

\linespread{1.3}\zihao{-4}\bigskip

This report studies the joint optimization of inspection intervals and tool replacement policies for an automated lathe process. Based on 100 recorded tool-failure lifetimes and given cost parameters, we formulate an expected-cost model that minimizes the average cost per qualified product. The model includes defective-product loss, inspection cost, fault-recovery cost, and preventive replacement cost. Misclassification loss is further incorporated in the second scenario.

The solution framework consists of lifetime distribution estimation, policy parameterization, and two-level optimization: first optimize inspection points for a fixed number of inspections, then search across different numbers of inspections to obtain the global best policy. Two quality-identification scenarios are compared, and improved strategies such as double-check confirmation are discussed.

The study shows that the lathe policy design can be treated as a constrained nonlinear optimization problem, and the optimal policy is highly sensitive to inspection reliability and cost coefficients. The proposed workflow is applicable to preventive maintenance and quality control in discrete manufacturing systems.

\bigskip
\noindent{\zihao{4}\heiti Keywords:}
tool lifetime; inspection policy; replacement policy; expected cost; nonlinear optimization
